\documentclass{article}
\usepackage{graphicx} % Required for inserting images

\title{Programations and Paradigms}
\author{Andrés López Gónzalez}
\date{August 2025}

\begin{document}

@Aprender sobre dificultad logaritmica, busqueda binaria, bubble sort
@''- Algoritmos en general
@Josh Camus
@Bit de sobreflujo, estudiar complemento 1 y 2

Que es complemento 2 (C2)?, como suma las cosas la maquina, y como puede haber un error en complemento A2 que desemboca negativos es (ini+fin)/2.
Existe una solucion pero provoca overflow.

La solucion se plantea como m = (i+f+i/i)/2 = i + (f_i)/2

De esta manera se arregla el overflow para todos los casos.
NO habian tantos datos desde el 1946 13feb que se creo la ENIAC primera computadora, por eso hasta casi 60 anos se dieron cuenta que habian problemas de overflow

white two key se soluciono de una forma similar, para entrar en contexto este fue **fill. Se soluciona la caida de la busqueda binaria con i + (f_i)/2.
complemento a 2 es la representacion de mathbb(Z) que garantiza. C2 es porque debemos representar un -n en algun momento, necesitamos al menos un bit adicional que llamamos bit de signo. Este regularmente es agregar un digito al inicio del numero.

por ej si 50 es 110010, puede extenderse como 0110010, un 0 mudo pero que garantiza positividad, -50 seria entonces 1110010.

Si existe C2 es intuitivo que existe un C1, tomas el numero positivo y tomas el complemento. Ya con el signo incluido si no no funciona, tomamos el numero que es positivo y solo lo invertimos.
1001101
+     1 = 1001110

Se lo debemos a leibniz y pascal con su primera calculadora, pues habia un problema con la resta, este era **fill




\end{document}
